% GENERATED FILE. Edit canonical lessons, then run npm run generate:books.
% canonical-source-hash: fnv1a64:618b8ec25da86844
% canonical-lessons: IT-C426-trolley, IT-C426-stiva, IT-C426-nastro, IT-C426-bilancia, IT-C426-imballare

\chapter{I bagagli}
\label{ch:it-i-bagagli}
\hlchaptermodality{426}

\begin{chapteropening}
\textbf{By the end of this chapter:} \emph{I can deal with luggage: the wheeled case, the hold, the baggage belt, the scales and packing things up.}

Complete the last lesson of chapter 426: I can deal with luggage: the wheeled case, the hold, the baggage belt, the scales and packing things up.
\end{chapteropening}

\section[il trolley]{il trolley — the wheeled suitcase}

\label{lesson:IT-C426-trolley}



\begin{quote}
Before the new one: say the Italian for the shuttle, then the Italian for the time zone.
\end{quote}

\subsection*{You'll want to know: il trolley}
\textbf{il trolley} — \textquotedblleft{}the wheeled suitcase\textquotedblright{}.

\begin{quote}
Porto solo un trolley piccolo.
\end{quote}

\begin{quote}
\emph{I am taking only a small wheeled case.}
\end{quote}

An English word with an Italian meaning: in English a trolley is a cart.

\begin{cousinweb}
\emph{trolley} was borrowed from English \emph{trolley} 'a small cart', and in Italian it came to mean a suitcase on wheels.
\end{cousinweb}

\begin{grammarlens}[title={What you've built}]
One more word for telling, explaining and describing at B1.
\end{grammarlens}

\subsection*{Your turn}
\begin{itemize}
\raggedright
  \item \emph{Say it:} \emph{il trolley}
  \item \emph{Say it:} \emph{il trolley}, once more
  \item \emph{Say it:} a sentence of your own with \emph{il trolley}
  \item \emph{Recall:} say the Italian for the shuttle, then the Italian for the time zone, then say \emph{il trolley} again
\end{itemize}

\begin{tcolorbox}[breakable,colback=teal!4,colframe=teal!35!black,title={Before you move on}]
What does il trolley mean? (\textquotedblleft{}The wheeled suitcase\textquotedblright{}.) Say it once more.
\end{tcolorbox}
\section[la stiva]{la stiva — the hold (of a ship or a plane)}

\label{lesson:IT-C426-stiva}



\begin{quote}
Before the new one: say the Italian for the time zone, then the Italian for the wheeled suitcase.
\end{quote}

\subsection*{You'll want to know: la stiva}
\textbf{la stiva} — \textquotedblleft{}the hold (of a ship or a plane)\textquotedblright{}.

\begin{quote}
Le valigie grandi vanno in stiva.
\end{quote}

\begin{quote}
\emph{The big suitcases go in the hold.}
\end{quote}

\emph{Il bagaglio da stiva} — checked luggage.

\begin{cousinweb}
\emph{stiva} is built on \emph{stivare} 'to stow, to load', which is inherited from Latin \emph{stīpāre} 'to press together' (EtymDB).
\end{cousinweb}

\begin{grammarlens}[title={What you've built}]
One more word for telling, explaining and describing at B1.
\end{grammarlens}

\subsection*{Your turn}
\begin{itemize}
\raggedright
  \item \emph{Say it:} \emph{la stiva}
  \item \emph{Say it:} \emph{la stiva}, once more
  \item \emph{Say it:} a sentence of your own with \emph{la stiva}
  \item \emph{Recall:} say the Italian for the time zone, then the Italian for the wheeled suitcase, then say \emph{la stiva} again
\end{itemize}

\begin{tcolorbox}[breakable,colback=teal!4,colframe=teal!35!black,title={Before you move on}]
What does la stiva mean? (\textquotedblleft{}The hold (of a ship or a plane)\textquotedblright{}.) Say it once more.
\end{tcolorbox}
\section[il nastro]{il nastro — the belt (for baggage); the ribbon}

\label{lesson:IT-C426-nastro}



\begin{quote}
Before the new one: say the Italian for the wheeled suitcase, then the Italian for the hold.
\end{quote}

\subsection*{You'll want to know: il nastro}
\textbf{il nastro} — \textquotedblleft{}the belt (for baggage); the ribbon\textquotedblright{}.

\begin{quote}
Aspettiamo le valigie al nastro sei.
\end{quote}

\begin{quote}
\emph{We wait for the suitcases at belt six.}
\end{quote}

\emph{Il nastro trasportatore} — conveyor belt; also a ribbon for hair or gifts.

\begin{cousinweb}
\emph{nastro} was borrowed from Gothic \emph{nastilō} 'strap, band' (EtymDB), a Germanic word.
\end{cousinweb}

\begin{grammarlens}[title={What you've built}]
One more word for telling, explaining and describing at B1.
\end{grammarlens}

\subsection*{Your turn}
\begin{itemize}
\raggedright
  \item \emph{Say it:} \emph{il nastro}
  \item \emph{Say it:} \emph{il nastro}, once more
  \item \emph{Say it:} a sentence of your own with \emph{il nastro}
  \item \emph{Recall:} say the Italian for the wheeled suitcase, then the Italian for the hold, then say \emph{il nastro} again
\end{itemize}

\begin{tcolorbox}[breakable,colback=teal!4,colframe=teal!35!black,title={Before you move on}]
What does il nastro mean? (\textquotedblleft{}The belt (for baggage); the ribbon\textquotedblright{}.) Say it once more.
\end{tcolorbox}
\section[la bilancia]{la bilancia — the scales}

\label{lesson:IT-C426-bilancia}



\begin{quote}
Before the new one: say the Italian for the hold, then the Italian for the belt.
\end{quote}

\subsection*{You'll want to know: la bilancia}
\textbf{la bilancia} — \textquotedblleft{}the scales\textquotedblright{}.

\begin{quote}
Metti la valigia sulla bilancia.
\end{quote}

\begin{quote}
\emph{Put the suitcase on the scales.}
\end{quote}

\emph{La bilancia pesapersone} — bathroom scales.

\begin{cousinweb}
\emph{bilancia} is inherited from Late Latin \emph{bilancia} 'balance', from \emph{bilanx} 'having two pans', built from \emph{bi-} and \emph{lanx} 'dish' (EtymDB).
\end{cousinweb}

\begin{grammarlens}[title={What you've built}]
One more word for telling, explaining and describing at B1.
\end{grammarlens}

\subsection*{Your turn}
\begin{itemize}
\raggedright
  \item \emph{Say it:} \emph{la bilancia}
  \item \emph{Say it:} \emph{la bilancia}, once more
  \item \emph{Say it:} a sentence of your own with \emph{la bilancia}
  \item \emph{Recall:} say the Italian for the hold, then the Italian for the belt, then say \emph{la bilancia} again
\end{itemize}

\begin{tcolorbox}[breakable,colback=teal!4,colframe=teal!35!black,title={Before you move on}]
What does la bilancia mean? (\textquotedblleft{}The scales\textquotedblright{}.) Say it once more.
\end{tcolorbox}
\section[imballare]{imballare — to pack, to wrap (for transport)}

\label{lesson:IT-C426-imballare}



\begin{quote}
Before the new one: say the Italian for the belt, then the Italian for the scales.
\end{quote}

\subsection*{You'll want to know: imballare}
\textbf{imballare} — \textquotedblleft{}to pack, to wrap (for transport)\textquotedblright{}.

\begin{quote}
Imballiamo bene i piatti per il viaggio.
\end{quote}

\begin{quote}
\emph{We pack the plates carefully for the journey.}
\end{quote}

\emph{L'imballaggio} — the packaging. For clothes in a case you say \emph{fare la valigia}.

\begin{cousinweb}
\emph{imballare} is built from \emph{in-} and \emph{balla} 'bale', which is usually traced to a Germanic word that came into Italian as a borrowing.
\end{cousinweb}

\begin{grammarlens}[title={What you've built}]
One more word for telling, explaining and describing at B1.
\end{grammarlens}

\subsection*{Your turn}
\begin{itemize}
\raggedright
  \item \emph{Say it:} \emph{imballare}
  \item \emph{Say it:} \emph{imballare}, once more
  \item \emph{Say it:} a sentence of your own with \emph{imballare}
  \item \emph{Recall:} say the Italian for the belt, then the Italian for the scales, then say \emph{imballare} again
\end{itemize}

\begin{tcolorbox}[breakable,colback=teal!4,colframe=teal!35!black,title={Before you move on}]
What does imballare mean? (\textquotedblleft{}To pack, to wrap (for transport)\textquotedblright{}.) Say it once more.
\end{tcolorbox}
